\documentclass[12pt,a4paper]{article}

\usepackage[utf8]{inputenc}
\usepackage[T1]{fontenc}
\usepackage[margin=2.5cm]{geometry}
\usepackage{hyperref}
\usepackage{microtype}
\usepackage{parskip}
\usepackage{lmodern}

\hypersetup{colorlinks=true, urlcolor=blue!60!black}

\pagestyle{empty}

\begin{document}

\begin{flushright}
Kundai Farai Sachikonye \\
Germany \\
\href{https://github.com/fullscreen-triangle}{github.com/fullscreen-triangle} \\
\texttt{kundai.sachikonye@bitspark.com} \\[6pt]
15 May 2026
\end{flushright}

\vspace{1em}

Dr.\ Hashem Mohammadian \\
Department of Medicine 3 \\
Universitätsklinikum Erlangen \\
Ulmenweg 18, 91054 Erlangen

\vspace{1.5em}

\textbf{Re: Bioinformatics PhD Position --- Computational Immunology (Job No.\ 12900)}

\vspace{1em}

Dear Dr.\ Mohammadian,

I am applying for the Bioinformatics PhD position in Computational Immunology.
My interest in this position is not general: the specific pathological focus of
your group --- autoimmune diseases, mesenchymal activation, inflammatory
arthritis, and fibrotic disease --- is the subject of theoretical and
computational work I have been developing independently over the past year,
and the methods your group employs are precisely those needed to test and
refine the predictions that work generates.

My paper \textit{On the Thermodynamic Consequences of Categorical Completion on
Biological Systems: Specificity through VDJ Recombination Categorical Filter
Cascades} proposes that adaptive immunity operates through categorical exclusion
rather than molecular complementarity. MHC molecules function as designed
ambiguous filters that selectively present peptides from proteins with high
categorical richness $R$ --- quantifying isoform diversity, conformational
entropy, and downstream connectivity --- while suppressing peptides from
low-$R$ self-proteins. The central empirical result is that MHC binding affinity
correlates with parent protein categorical richness independently of classical
sequence features: $r = 0.73$, $p < 10^{-12}$, $n = 2{,}847$ IEDB peptides for
MHC-I; $r = 0.68$, $p < 10^{-10}$ for MHC-II. Adding the categorical richness
term to NetMHCpan improves AUC from 0.68 to 0.81 ($p < 0.001$, DeLong test).
The framework predicts autoimmunity as the condition where self-proteins acquire
$R > 10^5$ and enter pathogen categorical space. I apply this explicitly to
collagen ($R \approx 10^{5.5}$), whose extensive isoform repertoire and dense
post-translational modification landscape push it across the threshold --- making
it a frequent autoimmune target by geometric necessity, not molecular accident.
This is the substrate of the pathology your group investigates.

A second paper, \textit{Syndrome: Categorical Resolution of Disease Trajectories},
formalises a computational framework for disease state representation using eight
cellular oscillator classes: protein folding, enzymatic catalysis, channel gating,
membrane potential, ATP synthesis, genetic expression, calcium signalling, and
circadian rhythm. Each class contributes a coherence index $\eta_i \in [0,1]$,
forming a disease vector $\mathbf{D} \in [0,1]^8$ that encodes the pathological
signature across biological scales. For inflammatory arthritis and mesenchymal
activation, the relevant components are immediate: elevated $D_E$ captures
metalloprotease and cytokine cascade dysregulation; elevated $D_{Ca}$ captures
the calcium-dependent fibroblast activation central to synovial pathology;
elevated $D_G$ captures the inflammatory gene expression signatures directly
measurable by single-cell RNA-sequencing. The framework computes disease
trajectories in $\mathcal{O}(N \log N)$ and selects optimal therapeutic
intervention targets as $\arg\max_i w_i D_i$ --- a sparse linear programme
with a computable a priori side-effect bound.

A third paper, \textit{On Disease Epistemology in Blind Receiver: Cellular
Ignorance, Group Blindness, and Self-Consistency as the Unique Basis of Disease
Detection}, proves that template-based diagnostics are epistemologically
impossible at every biological scale and derives the unique viable criterion:
loop-holonomy self-consistency of the biochemical circuit. A cell is healthy if
and only if the composed constraint map around every reaction cycle returns the
identity ($H_\ell = \mathrm{Id}\ \forall\ell$); disease is the condition
$H_\ell \neq \mathrm{Id}$ for some $\ell$. Numerically, the holonomy test
achieves AUC $= 1.00$ against AUC $= 0.896$ for template comparison across
25 validation experiments, all of which pass. This offers a fundamentally
different diagnostic strategy from current biomarker panels --- one that detects
circuit-level inconsistency that is locally invisible to individual-node
measurement, which is precisely the kind of inconsistency characteristic of
chronic autoimmune processes.

The computational infrastructure behind this theoretical work is directly
applicable to the experimental modalities your group uses.
\href{https://github.com/fullscreen-triangle/gospel}{gospel} processes
whole-genome VCF files alongside RNA-seq expression matrices at $10^6$
variants/second, validated against gnomAD v3.1.2 and UK Biobank;
the genetic expression oscillator in the Syndrome framework takes its input
from exactly this kind of data.
\href{https://github.com/fullscreen-triangle/lavoisier}{lavoisier} is a
mass-spectrometry metabolomics pipeline achieving 98.9\% feature extraction
accuracy against NIST reference libraries, with distributed computing via Ray
and PyTorch-based neural networks.
\href{https://github.com/fullscreen-triangle/helicopter}{helicopter} addresses
multi-modal microscopy image analysis across fluorescence, brightfield, and
spectral modalities --- directly applicable to imaging mass cytometry data.
I have also built deployable web tools for sequence homology search
(\href{https://vivid-symbolism.vercel.app/search}{vivid-symbolism.vercel.app/search}),
multi-modal biological data observation
(\href{https://hieronymus-omega.vercel.app/observe}{hieronymus-omega.vercel.app/observe}),
and protein analysis --- folding, trajectory, catalysis, and dynamics via Kuramoto
oscillator synchronization in S-entropy space
(\href{https://shakespear-nine.vercel.app/instrument}{shakespear-nine.vercel.app/instrument})
--- all available in the browser without installation, demonstrating that I build
accessible, reproducible research tools rather than scripts.

My academic background is in computational lipidomics at the Technische
Universität München, where I trained in the statistical and multi-omics
methods that underpin all of this work. R and Python are both standard tools
from that training. I am not proposing to import my framework wholesale into
your group's programme --- that is not what a doctoral position is for. I
am proposing that a candidate who has independently derived and empirically
validated a mechanistic theory of why collagen and mesenchymal proteins become
autoimmune targets, and built the computational infrastructure to analyse the
data your group generates, will contribute productively to that programme from
the first week.

I am available from July 2026, fully authorised to work in Germany
(Aufenthaltserlaubnis), and would welcome the opportunity to discuss
this work with you directly.

\vspace{1.5em}
Yours sincerely,

\vspace{2.5em}
\textbf{Kundai Farai Sachikonye}

\end{document}
